\section{Conclusiones Generales}

A partir del análisis sistemático y secuencial de las simulaciones fundamentales y los datos de dispersión molecular, se extraen las siguientes conclusiones integrales que abarcan todos los fenómenos estudiados:

\textbf{Efecto Túnel y Validación Numérica:} La simulación Crank-Nicolson 1D reproduce con alta precisión la transmisión para una barrera rectangular con $V_0 > E$. El barrido paramétrico confirma que $T$ decrece $V_0$ y $w$, en concordancia con la  aproximación WKB. Por otro lado, la conservación de probabilidad $T+R+P_\text{barrera}\approx 1$ durante toda la evolución atestigua la estabilidad y fiabilidad del esquema implícito empleado.
 
\textbf{Dinámica Cuántica en la Doble Rendija:} La evolución espaciotemporal de la función de onda $\psi(x,y)$ demuestra fenomenológicamente el principio de superposición cuántica. La división topológica del paquete de ondas incidente genera dos frentes de onda coherentes que interactúan, formando un patrón de interferencia determinista. Las transiciones de fase ($\arg(\psi) \in [-\pi, \pi]$) evidencian que el sistema conserva la coherencia de fase global durante la propagación, validando de forma visual las predicciones de la ecuación de Schrödinger dependiente del tiempo para la difracción de materia pura.


\textbf{Precisión de la Simulación de Monte Carlo:} El análisis cuantitativo del perfil transversal de intensidades ($I/I_0$) constata una convergencia excepcional entre la detección acumulada de partículas y el modelo teórico policromático promediado. La morfología del pico central (orden cero) y las oscilaciones secundarias adyacentes certifican que el factor de estructura asociado a la periodicidad atómica del HBC y el factor de forma de las rendijas fueron parametrizados correctamente, evidenciando un evento de difracción macromolecular coherente.


La concatenación de estos resultados establece un puente directo entre los principios fundamentales de la interferencia (doble rendija) y sus aplicaciones cristalográficas avanzadas (nanografeno). Queda demostrado empírica y teóricamente que la difracción de materia coherente constituye un método analítico riguroso y no destructivo, capaz de trasladar las propiedades geométricas microscópicas al espacio recíproco para la caracterización precisa de materiales bidimensionales complejos.



